\section*{Overview}

Primitive roots and discrete logarithms are two sides of the same idea: if a multiplicative modular group is cyclic,
then every non-zero element can be represented as a power of one generator. That turns multiplication into addition on
exponents.

\section*{When to Use It}

Use these ideas when:

\begin{itemize}
  \item the problem has equations like \(a^x \equiv b \pmod m\),
  \item you need to enumerate powers by generator order,
  \item a multiplicative subgroup needs to be indexed explicitly.
\end{itemize}

\section*{Core Idea}

For a prime modulus \(p\), the non-zero residues form a cyclic group of size \(p-1\). If \(g\) is a primitive root
modulo \(p\), then every non-zero residue is \(g^k\) for some \(k\).

The discrete logarithm problem asks to recover \(k\) from
\[
g^k \equiv h \pmod p.
\]
The standard contest algorithm is baby-step giant-step.

\section*{Key Insight}

Once a generator exists, multiplicative equations become additive equations on exponents modulo the group order. The
hard part shifts from algebra to fast lookup of powers.

\section*{Worked Problem}

\paragraph{Problem.}
For a prime \(p\), given \(g\) and \(h\), find the smallest non-negative \(x\) such that
\[
g^x \equiv h \pmod p,
\]
or report that no such \(x\) exists.

\paragraph{Why these tools fit.}
If \(g\) generates the relevant subgroup, the problem is exactly a discrete logarithm. Baby-step giant-step rewrites
\(x = iq + j\) and matches
\[
g^{iq} \equiv h g^{-j}.
\]

\section*{Correctness Intuition}

Baby-step giant-step covers every exponent by splitting it into a large-block index \(i\) and a small offset \(j\). A
hash table over one side and a linear scan over the other finds the match in \(O(\sqrt{m})\) time.

\section*{Complexity Analysis}

\begin{itemize}
  \item primitive root modulo prime \(p\): factor \(p-1\), then test candidates in roughly \(O(\tau(p-1)\log p)\),
  \item baby-step giant-step: \(O(\sqrt{m}\log m)\) time and \(O(\sqrt{m})\) memory.
\end{itemize}

\section*{Implementation}

The code includes:

\begin{itemize}
  \item \texttt{primitive\_root\_prime(p)} for prime moduli,
  \item \texttt{discrete\_log\_prime(g, h, p)} with baby-step giant-step.
\end{itemize}

\section*{Common Pitfalls}

\begin{itemize}
  \item Forgetting the primitive-root existence conditions for general composite moduli.
  \item Calling BSGS when \(g\) and \(p\) are not coprime in the implemented variant.
  \item Ignoring subgroup issues; not every \(h\) is reachable from every \(g\).
\end{itemize}

\section*{Variants / Extensions}

\begin{itemize}
  \item General discrete log with gcd handling.
  \item Using exponent indices to solve multiplicative congruences.
  \item Primitive roots for NTT-friendly moduli and subgroup generators.
\end{itemize}

\section*{Practice Problems}

\begin{itemize}
  \item Solve \(a^x \equiv b \pmod p\).
  \item Find a generator of \((\mathbb{Z}/p\mathbb{Z})^\times\).
  \item Transform multiplicative constraints into additive ones over exponents.
\end{itemize}

\section*{References}

\begin{itemize}
  \item \href{https://cp-algorithms.com/algebra/discrete-log.html}{cp-algorithms: Discrete Logarithm}.
  \item \href{https://cp-algorithms.com/algebra/primitive-root.html}{cp-algorithms: Primitive Root}.
  \item \href{https://codeforces.com/catalog}{Codeforces Catalog}.
\end{itemize}
